\documentclass[11pt,letterpaper]{article}

% ---------- Idioma y codificación ----------
\usepackage[utf8]{inputenc}
\usepackage[T1]{fontenc}
\usepackage[spanish,es-tabla]{babel}

% ---------- Paquetes generales ----------
\usepackage[margin=2.5cm]{geometry}
\usepackage{xcolor}
\usepackage{enumitem}
\usepackage{fancyhdr}
\usepackage{titlesec}
\usepackage{tcolorbox}
\tcbuselibrary{skins,breakable}
\usepackage{listings}
\usepackage{graphicx}
\usepackage{array}
\usepackage{longtable}
\usepackage{hyperref}
\usepackage{xurl}
\usepackage{seqsplit}

% ---------- Colores del curso ----------
\definecolor{brandblue}{HTML}{1F4E78}
\definecolor{brandgreen}{HTML}{1DB954}
\definecolor{brandorange}{HTML}{E07B00}
\definecolor{lightgray}{HTML}{F4F6F8}
\definecolor{codebg}{HTML}{F5F5F5}

% ---------- Encabezado y pie de página ----------
\pagestyle{fancy}
\fancyhf{}
\renewcommand{\headrulewidth}{0.4pt}
\fancyhead[L]{\small Taller Guiado --- Analítica de Datos}
\fancyhead[R]{\small Cierre de capítulos en Orange}
\fancyfoot[C]{\small\thepage}

% ---------- Estilo de secciones ----------
\titleformat{\section}
  {\normalfont\Large\bfseries\color{brandblue}}
  {\thesection}{0.6em}{}
\titleformat{\subsection}
  {\normalfont\large\bfseries\color{brandblue}}
  {\thesubsection}{0.6em}{}

% ---------- Cajas de actividades, preguntas y notas ----------
\newtcolorbox{actividad}[1][]{
  breakable, enhanced, colback=lightgray, colframe=brandblue,
  boxrule=0.8pt, arc=2pt, left=6pt, right=6pt, top=5pt, bottom=5pt,
  title={\faPencil\ Actividad}, fonttitle=\bfseries\color{white},
  colbacktitle=brandblue, #1
}
\newtcolorbox{pregunta}[1][]{
  breakable, enhanced, colback=white, colframe=brandorange,
  boxrule=0.8pt, arc=2pt, left=6pt, right=6pt, top=5pt, bottom=5pt,
  title={\faPencil\ Pregunta de concepto}, fonttitle=\bfseries\color{white},
  colbacktitle=brandorange, #1
}
\newtcolorbox{nota}[1][]{
  breakable, enhanced, colback=white, colframe=brandgreen,
  boxrule=0.8pt, arc=2pt, left=6pt, right=6pt, top=5pt, bottom=5pt,
  title={Nota}, fonttitle=\bfseries\color{white},
  colbacktitle=brandgreen, #1
}

% ---------- Código / comandos ----------
\lstset{
  basicstyle=\ttfamily\small,
  backgroundcolor=\color{codebg},
  breaklines=true,
  frame=single,
  framerule=0.3pt,
  rulecolor=\color{gray!50},
  columns=fullflexible,
  keepspaces=true,
}

% ---------- Metadatos del PDF ----------
\hypersetup{
  pdftitle={Taller Guiado - Cierre de capitulos en Orange - Analitica de Datos},
  pdfauthor={Diego Fabián Collazos Huertas},
  pdfsubject={Curso Analítica de Datos - Taller guiado},
  colorlinks=true,
  linkcolor=brandblue,
  urlcolor=brandblue,
  citecolor=brandblue,
}

% Sustituto simple para el icono de lápiz (evita depender de fontawesome)
\newcommand{\faPencil}{\raisebox{-0.1em}{\textbf{\Large\textbullet}}}

\title{\vspace{-2cm}}
\date{}

\begin{document}

% ============================================================
% PORTADA
% ============================================================
\begin{titlepage}
\centering
\vspace*{1.5cm}

{\Large Universidad Nacional de Colombia --- Sede Manizales}\\[0.2cm]
{\large Facultad de Ingeniería y Arquitectura}\\[0.2cm]
{\large Departamento de Ingeniería Eléctrica, Electrónica y Computación}\\[1.2cm]

\rule{\linewidth}{1pt}\\[0.5cm]
{\huge\bfseries\color{brandblue} Taller Guiado en Orange}\\[0.4cm]
{\Large Cierre de capítulos: un repaso de extremo a extremo}\\[0.5cm]
\rule{\linewidth}{1pt}\\[1cm]

{\large \textbf{Curso:} Analítica de Datos (Aprendizaje de máquina) --- Código 4200729}\\[0.3cm]
{\large \textbf{Docente:} Diego Fabián Collazos Huertas}\\[0.3cm]
{\large \textbf{Modalidad:} guiado en clase, individual o en parejas}\\[0.3cm]
{\large \textbf{Dataset:} \texttt{ICFES\_2020-2\_Eje\_Cafetero.csv}}\\[1.5cm]

\vfill

{\small Este taller se desarrolla \textbf{en el horario de clase}, guiado en vivo por el docente. El archivo \texttt{.ows} resultante, con las respuestas a las preguntas de concepto escritas sobre el lienzo, se entrega según se indique en clase.}

\end{titlepage}

% ============================================================
% OBJETIVOS
% ============================================================
\section{Objetivos}

\begin{itemize}[leftmargin=1.2cm]
  \item Repasar, de extremo a extremo y sobre un \textbf{único problema real}, (casi) todas las técnicas de análisis exploratorio, selección de variables y agrupamiento vistas en el curso, usando \textbf{Orange Data Mining}.
  \item Practicar la formulación de preguntas sobre un conjunto de datos y la elección de la técnica adecuada para responderlas.
  \item Hacer un \textbf{sanity check} de los datos antes de sacar conclusiones, identificando valores extraños, atípicos y faltantes.
  \item Distinguir cuándo una diferencia observada es \textbf{estadísticamente significativa} y cuándo una correlación no implica causalidad.
\end{itemize}

\begin{nota}
Este enfoque \textbf{no es el único posible}: dependiendo de los datos y de las preguntas que queramos hacerles, las técnicas, los análisis y las conclusiones pueden cambiar. Lo importante de este taller es la \textbf{forma de pensar} el problema, no memorizar una receta fija.
\end{nota}

% ============================================================
% DATOS DE ORIGEN
% ============================================================
\section{Datos de origen}

Trabajaremos con un extracto de la base de datos de las \textbf{Pruebas Saber 11, segundo semestre de 2020} (origen: \href{https://www.datos.gov.co/Educaci-n/Resultados-nicos-Saber-11/kgxf-xxbg}{Saber 11° 2020-2 --- Datos Abiertos Colombia}). Se le hizo un depurado previo:

\begin{itemize}[leftmargin=1.2cm]
  \item Se seleccionaron solo registros del \textbf{Eje Cafetero} (Risaralda, Caldas, Quindío).
  \item Se eliminaron variables irrelevantes (códigos SNIES de colegios y variables similares).
  \item Se hizo un submuestreo al 40\% para simplificar el procesamiento.
\end{itemize}

El archivo resultante, \texttt{\seqsplit{ICFES\_2020-2\_Eje\_Cafetero.csv}}, queda con \textbf{24999 registros} y \textbf{61 variables} (40 categóricas, 20 numéricas, 1 metadato de texto): puntajes y percentiles por competencia (lectura crítica, matemáticas, ciencias naturales, sociales y ciudadanas, inglés, y el puntaje global), variables socioeconómicas y del entorno familiar del estudiante, y variables del colegio.

% ============================================================
% PARTE 1
% ============================================================
\section{Parte 1 --- Carga y lectura general de los datos}

\begin{enumerate}[leftmargin=1.2cm]
  \item Abre Orange y crea un workflow nuevo. Arrastra el widget \textbf{File} y carga \texttt{ICFES\_2020-2\_Eje\_Cafetero.csv}. Confirma en el panel del widget que quedan \textbf{24999 filas, 40 categóricas, 20 numéricas, 1 meta}.
  \item Agrega el widget \textbf{Feature Statistics} (categoría \textit{Data}) conectado a \textit{File} y revisa, variable por variable, la media, moda, mediana, dispersión, mínimo y máximo.
  \item Fíjate en las tres primeras variables de código (\texttt{ESTU\_COD\_RESIDE\_DEPTO}, \texttt{ESTU\_MCPIO\_RESIDE}, \texttt{ESTU\_COD\_RESIDE\_MCPIO}): aunque Orange las detecta como \textbf{Numéricas} por estar hechas de dígitos, no tiene sentido promediarlas --- son códigos, no cantidades. Agrega un widget \textbf{Edit Domain} conectado a \textit{File} y cambia su tipo a \textbf{Categórica}.
\end{enumerate}

\begin{pregunta}
¿Por qué un código numérico como \texttt{ESTU\_COD\_RESIDE\_MCPIO} no debería tratarse como variable numérica, aunque esté compuesto solo por dígitos? Den un ejemplo concreto de una conclusión sin sentido que se obtendría si se calculara su \textbf{media} o se le aplicara una \textbf{correlación}.
\end{pregunta}

% ============================================================
% PARTE 2
% ============================================================
\section{Parte 2 --- Sanity check: ¿los datos son lo que deberían ser?}

Antes de analizar, conviene verificar que no haya datos mal tomados o en un formato incorrecto --- a esto se le llama \textbf{sanity check}.

\begin{enumerate}[leftmargin=1.2cm]
  \item En \textit{Feature Statistics}, observa la variable \texttt{ESTU\_FECHANACIMIENTO}: su \textbf{dispersión} es de aproximadamente 74 años, con un mínimo en 1946 y un máximo en 2020.
  \item Agrega el widget \textbf{Select Rows} conectado a \textit{File} con la condición \texttt{ESTU\_FECHANACIMIENTO is at most 1970-01-01}. Conecta la salida a un \textbf{Data Table} e inspecciona esos registros.
  \item Conecta ahora un widget \textbf{Distributions} a la salida de \textit{Select Rows} y mira la variable \texttt{ESTU\_MCPIO\_RESIDE}. Notarás un eje X muy ruidoso, con muchas categorías sin barra: Orange sigue mostrando \textbf{todos} los municipios posibles del dataset completo, aunque esta selección casi no tenga datos en la mayoría. En el panel de \textit{Select Rows}, activa \textbf{Remove unused features} y observa cómo se limpia la visualización.
\end{enumerate}

\begin{pregunta}
¿Por qué es importante hacer un sanity check \textbf{antes} de interpretar los resultados de un análisis exploratorio? En este dataset hay personas que, según la fecha de nacimiento, tenían 74 años al presentar la prueba, y otras que nacieron el mismo año en que la presentaron (2020). ¿Cuál de las dos situaciones les parece más plausible como dato real y cuál como posible error de captura? Justifiquen.
\end{pregunta}

% ============================================================
% PARTE 3
% ============================================================
\section{Parte 3 --- Visualización de variables categóricas}

\subsection{3.1 Distributions}

\begin{enumerate}[leftmargin=1.2cm]
  \item Conecta un widget \textbf{Distributions} a \textit{File} (sin el filtro de la Parte 2) y selecciona \texttt{ESTU\_GENERO}.
  \item Usa \textbf{Split by} para condicionar la distribución de \texttt{FAMI\_ESTRATOVIVIENDA} según \texttt{ESTU\_GENERO}, y explora las opciones \textit{Sort categories by frequency}, \textit{Stack columns} y \textit{Show probabilities}.
\end{enumerate}

\begin{pregunta}
¿Qué diferencia hay entre una distribución \textbf{marginal} (una sola variable, como el histograma de \texttt{ESTU\_GENERO}) y una distribución \textbf{condicional} (\textit{split by}, como \texttt{FAMI\_ESTRATOVIVIENDA} dividida por género)? ¿Qué pregunta puede responder la segunda que la primera no puede?
\end{pregunta}

\subsection{3.2 Mosaic}

\begin{enumerate}[leftmargin=1.2cm]
  \item \textit{Distributions} queda limitado a dos variables (una principal y un \textit{split by}). Agrega el widget \textbf{Mosaic} conectado a \textit{File} para visualizar más variables a la vez (hasta 5): selecciona \texttt{ESTU\_GENERO} y \texttt{ESTU\_MCPIO\_RESIDE}, con color interior \texttt{ESTU\_GENERO}.
  \item Activa \textbf{Compare with total} y observa cómo cambia la lectura del gráfico.
  \item Agrega una tercera variable (por ejemplo \texttt{ESTU\_FECHANACIMIENTO} discretizada, o cualquier otra categórica) y luego hasta una cuarta o quinta. Prueba el botón \textbf{Find Informative Mosaics} --- \textbf{con pocas variables y datos reducidos}, ya que puede tardar mucho con el dataset completo.
\end{enumerate}

\begin{pregunta}
\textit{Mosaic} permite visualizar hasta 5 variables categóricas a la vez, mientras que \textit{Distributions} solo permite 2. ¿Qué ventaja ofrece poder ver más variables al tiempo? ¿Qué desventaja, en términos de qué tan fácil es \textbf{interpretar} el gráfico, aparece a medida que agregan más variables?
\end{pregunta}

% ============================================================
% PARTE 4
% ============================================================
\section{Parte 4 --- Variables numéricas: atípicos y faltantes}

\subsection{4.1 Box Plot}

\begin{enumerate}[leftmargin=1.2cm]
  \item Agrega un widget \textbf{Select Rows} conectado a \textit{File} para quedarte solo con estudiantes nacidos entre 2000 y 2004.
  \item Conecta un \textbf{Box Plot} a esa salida: variable \texttt{PUNT\_GLOBAL}, \textbf{Subgroups} = \texttt{\seqsplit{ESTU\_COD\_RESIDE\_DEPTO}}, con \textbf{Compare means} activado. Observa el valor de \textbf{ANOVA} que reporta el widget.
  \item Prueba también un \textit{Box Plot} con una variable \textbf{categórica} como variable principal (por ejemplo \texttt{COLE\_GENERO}) agrupada por otra categórica (\texttt{ESTU\_COD\_RESIDE\_DEPTO}): el widget reporta entonces un estadístico $\chi^2$ en vez de ANOVA.
\end{enumerate}

\begin{pregunta}
El \textit{Box Plot} de \texttt{PUNT\_GLOBAL} por departamento reporta un ANOVA con cierto valor $p$. Si ese valor $p$ es mayor a 0.05, ¿pueden concluir que hay una diferencia \textbf{estadísticamente significativa} entre los departamentos al 95\% de confianza? ¿Qué significa, en términos simples, el valor $p$ de una prueba de significancia?
\end{pregunta}

\subsection{4.2 Outliers y valores faltantes}

\begin{enumerate}[leftmargin=1.2cm]
  \item Conecta un \textbf{Distributions} a \textit{File} con la variable \texttt{PUNT\_GLOBAL}. Observa las colas de la distribución: hay valores muy por debajo de 100 que no deberían ser posibles en una escala que llega a 500.
  \item Agrega un \textbf{Select Rows} con la condición \texttt{PUNT\_GLOBAL is at most 100} y revisa esos registros en un \textit{Data Table}. Decide si te parecen datos reales o errores de captura.
  \item Agrega el widget \textbf{Outliers} (categoría \textit{Unsupervised}) conectado a \textit{File} para detectar valores atípicos considerando \textbf{todas} las variables numéricas a la vez, no solo una.
  \item Con un \textbf{Data Table} conectado a \textit{File}, identifica qué variables tienen valores faltantes --- notarás que están relacionadas con el entorno socioeconómico, no con los puntajes. Agrega un widget \textbf{Impute} para tratarlos.
\end{enumerate}

\begin{nota}
El widget \textit{Outliers} es \textbf{pesado}: puede ser lento o bloquear Orange si se usa con muchos datos y muchas variables. Redúcelo primero (\textit{Select Columns} y/o una muestra con \textit{Data Sampler}) antes de usarlo sobre el dataset completo.
\end{nota}

\begin{pregunta}
¿Cuál es la diferencia entre un \textbf{valor atípico legítimo} (una observación real, pero poco común) y un \textbf{error de captura de datos}? En el caso de los puntajes globales por debajo de 100, ¿qué evidencia usarían para decidir a cuál de los dos casos corresponde, en vez de simplemente eliminarlos sin revisar?
\end{pregunta}

% ============================================================
% PARTE 5
% ============================================================
\section{Parte 5 --- Formulación de preguntas y significancia estadística}

Una vez hecha la lectura general de los datos, el análisis exploratorio se vuelve más potente cuando parte de \textbf{preguntas concretas}. Algunos ejemplos:

\begin{itemize}[leftmargin=1.2cm]
  \item ¿Los estudiantes que sacan puntajes altos comparten características socioeconómicas similares?
  \item ¿Tener buenos hábitos de lectura es indicador de que me va a ir bien en el ICFES?
  \item ¿Estudiar en colegio privado es una variable que tiene correlación con sacar buenos puntajes?
  \item ¿Las personas con mejores puntajes en competencias ciudadanas vienen de colegios públicos o privados?
  \item ¿Hay relación entre el estrato y los resultados en competencias ciudadanas?
  \item ¿La alimentación familiar podría tener que ver con los resultados de las pruebas?
\end{itemize}

Tomemos la pregunta \textbf{¿Estudiar en colegio privado tiene relación con sacar buenos puntajes?} \texttt{COLE\_NATURALEZA} es \textbf{categórica} y \texttt{PUNT\_GLOBAL} es \textbf{numérica}, así que \textbf{no podemos calcular directamente un coeficiente de correlación} entre ellas.

\begin{enumerate}[leftmargin=1.2cm]
  \item Conecta un \textbf{Distributions} a \textit{File}: variable \texttt{PUNT\_GLOBAL}, \textbf{Split by} = \texttt{COLE\_NATURALEZA}. Observa cómo se desplazan las dos distribuciones condicionales una respecto a la otra.
  \item Conecta un \textbf{Box Plot}: variable \texttt{PUNT\_GLOBAL}, \textbf{Subgroups} = \texttt{COLE\_NATURALEZA}, con \textbf{Compare means}. Esta vez el widget sí reporta una prueba de significancia (t de Student) entre las dos medias.
  \item Repite el mismo \textit{Box Plot} para la pregunta sobre competencias ciudadanas (\texttt{\seqsplit{PUNT\_SOCIALES\_CIUDADANAS}} por \texttt{COLE\_NATURALEZA}).
\end{enumerate}

\begin{pregunta}
Podríamos codificar (\textit{Continuize}) la variable categórica \texttt{COLE\_NATURALEZA} como 0/1 y calcular su correlación de Pearson con \texttt{PUNT\_GLOBAL}. ¿Por qué el número resultante no significaría realmente lo que esperamos? Y aunque el \textit{Box Plot} muestre una diferencia de medias \textbf{estadísticamente significativa} entre colegios oficiales y no oficiales, ¿pueden concluir que estudiar en colegio privado \textbf{causa} mejores puntajes? ¿Qué otras explicaciones alternativas se les ocurren?
\end{pregunta}

% ============================================================
% PARTE 6
% ============================================================
\section{Parte 6 --- Selección de variables y vecinos}

\subsection{6.1 Selección de variables con Rank}

Queremos identificar qué variables socioeconómicas distinguen mejor entre puntajes \textbf{altos}, \textbf{medios} y \textbf{bajos}. Para eso necesitamos una variable \textbf{Target} categórica.

\begin{enumerate}[leftmargin=1.2cm]
  \item Agrega el widget \textbf{Discretize} conectado a \textit{File}, y discretiza \texttt{PUNT\_GLOBAL} en tres categorías (por ejemplo, \textit{Equal frequency} con 3 intervalos: altos/medios/bajos).
  \item Con \textbf{Select Columns}, separa por un lado todas las variables \textbf{menos} \texttt{PUNT\_GLOBAL}, y por otro solo \texttt{PUNT\_GLOBAL} discretizada. Une ambas de nuevo con \textbf{Merge Data}.
  \item En el widget resultante (o con \textbf{Edit Domain}), marca la versión discretizada de \texttt{PUNT\_GLOBAL} como \textbf{Target}.
  \item Conecta un widget \textbf{Rank} y revisa qué variables socioeconómicas quedan mejor rankeadas para distinguir altos/medios/bajos.
\end{enumerate}

\begin{pregunta}
¿Por qué fue necesario discretizar \texttt{PUNT\_GLOBAL} (convertirla de numérica a categórica) antes de poder usarla como variable \textbf{Target} para el widget \textit{Rank}? ¿Qué le pasaría a Rank si le diéramos directamente la variable numérica como \textit{Target}?
\end{pregunta}

\subsection{6.2 Distancias y vecinos}

\begin{enumerate}[leftmargin=1.2cm]
  \item Con un \textbf{Select Rows}, separa el registro (o los registros) con el \textbf{mayor} \texttt{PUNT\_GLOBAL} de todo el dataset.
  \item Conecta un widget \textbf{Neighbors}: como \textit{Data} usa el resto del dataset, y como \textit{Reference} el registro de mayor puntaje. Elige la métrica \textbf{Jaccard} (apropiada para variables categóricas binarias) y revisa qué estudiantes resultan más ``cercanos''.
  \item Mira qué características socioeconómicas comparten esos vecinos con el estudiante de referencia.
\end{enumerate}

\begin{pregunta}
¿Por qué la distancia \textbf{Jaccard} es más apropiada que la distancia \textbf{Euclídea} para comparar variables categóricas binarias (como las de Sí/No: ¿tiene internet?, ¿tiene computador?, etc.)? (Pista: piensen en qué significaría restar numéricamente dos categorías como ``Sí'' y ``No''.)
\end{pregunta}

% ============================================================
% PARTE 7
% ============================================================
\section{Parte 7 --- Agrupamiento no supervisado}

\begin{enumerate}[leftmargin=1.2cm]
  \item \textbf{Con variables numéricas}: agrega \textbf{Select Columns} para quedarte solo con las variables numéricas de puntaje, luego \textbf{k-Means} (prueba varios $K$ y quédate con el mejor según el coeficiente \textbf{Silhouette}). Visualiza el resultado en un \textbf{Scatter Plot} coloreado por \texttt{Cluster}.
  \item \textbf{Con variables categóricas}: repite el ejercicio usando solo variables socioeconómicas (categóricas). Recuerda que \textit{k-Means} en Orange preprocesa automáticamente las categóricas (\textit{continuize} con una columna por valor) si no le conectas ya datos preprocesados.
  \item Un \textit{Scatter Plot} directo sobre dos variables categóricas es difícil de interpretar. Agrega un widget \textbf{MDS} (conectado a través de \textbf{Distances}) para proyectar esas variables a 2 dimensiones y visualizar los clústeres encontrados.
  \item Conecta un \textbf{Distributions} a la salida del agrupamiento categórico: variable \texttt{PUNT\_GLOBAL}, \textbf{Split by} = \texttt{Cluster}. Observa si los dos clústeres (formados \textbf{solo} con variables socioeconómicas, sin usar el puntaje) se relacionan con el puntaje global.
\end{enumerate}

\begin{pregunta}[title={\faPencil\ Pregunta de concepto final}]
Al agrupar usando \textbf{solo} las variables numéricas de puntaje, era esperable que \textit{k-Means} separara ``puntajes altos'' de ``puntajes bajos'' --- esas son justo las variables que usamos para agrupar. Pero al agrupar usando \textbf{solo} variables socioeconómicas (categóricas, sin usar el puntaje) y luego comparar esos clústeres con \texttt{PUNT\_GLOBAL}, ¿qué encontraron? ¿Qué les sugiere esto sobre la relación entre el contexto socioeconómico y el desempeño en las pruebas Saber 11? Y de nuevo: ¿esto es evidencia de que el contexto socioeconómico \textbf{causa} el desempeño, o solo de que están \textbf{correlacionados}?
\end{pregunta}

% ============================================================
% ENTREGABLES
% ============================================================
\section{Entregables}

\begin{enumerate}[leftmargin=1.2cm]
  \item El archivo \textbf{\texttt{.ows}} del workflow completo armado durante la sesión, con las respuestas a las \textbf{8 preguntas de concepto} (\faPencil) escritas directamente sobre el lienzo, usando la herramienta de texto de Orange (clic derecho sobre un espacio vacío $\to$ \textit{Text}), junto a la parte correspondiente del workflow.
  \item Sube el archivo \texttt{.ows} desde la página de entregas del curso, seleccionando la categoría \textbf{Otra} y especificando ``Taller guiado --- Cierre de capítulos'' --- o según indique el docente en clase.
\end{enumerate}

\begin{nota}
Este es un taller \textbf{guiado en clase}: no hay un procedimiento único ni una única forma correcta de armar el workflow. Lo importante es que el flujo refleje las decisiones que se discutieron en cada parte y que las respuestas a las preguntas de concepto estén justificadas con lo que efectivamente se observó en Orange.
\end{nota}

\end{document}
